\documentclass[11pt]{letter}

\usepackage[margin=0.75in]{geometry}
\usepackage{setspace}
\usepackage{hyperref}
\usepackage{xcolor}
\usepackage{enumitem}
\usepackage{parskip}

\hypersetup{
    colorlinks=true,
    urlcolor=black,
    linkcolor=black
}

\setlength{\parindent}{0pt}
\setstretch{1.05}
\date{}
\begin{document}

\begin{letter}{}

\opening{Dear Hiring Committee,}

I saw the job post on \textbf{Jobright} and thought I'd be a good fit for it.

I am a Master's student in Computer Science at the University of Cincinnati, graduating December 2026, looking for a role as a \textbf{Software Engineer 1, Cloud Infrastructure} at \textbf{Abnormal AI}. With two years of production cloud infrastructure experience at Siemens Industry Software, I am excited about the opportunity to contribute to \textbf{Abnormal AI} and support the \textbf{Platform and Infrastructure team}.

\textbf{Here are reasons that make me a great fit for the role:}

\begin{itemize}[leftmargin=*, itemsep=0.45em]
    \item \textbf{Infrastructure as code for reliable cloud systems:} At Siemens, I built Terraform and Ansible automation that standardized provisioning across 40+ production environments. As an intern, I cut deployment time by 30\% with reusable Terraform modules on AWS and built a Go and Terratest suite that caught 95\% of misconfigurations before production.

    \item \textbf{Python, AWS, and Kubernetes in production:} I automated Teamcenter X infrastructure workflows with Python (Boto3), replacing legacy Ruby scripts and cutting deployment time by 40\%. I also migrated 40+ production environments to Kubernetes using ArgoCD, Helm, GitLab CI/CD, and HashiCorp Vault.

    \item \textbf{Production reliability, cost, and less operational toil:} I executed 44 production environment upgrades in 4 months, diagnosing live issues with CloudWatch and Datadog. I designed a Blue-Green deployment framework that cut upgrade downtime by 50\%, reduced monthly AWS spend by 40\% (\$15K to \$9K), and built a self service platform adopted by 15 engineers that cut environment management time by up to 90\%.

    \item \textbf{Mature engineering practices and teamwork:} As Quality Coach and SonarQube Specialist, I raised automated code coverage from 35\% to 78\% and mentored 5 engineers through onboarding and standardized practices.

    \item \textbf{Judgment when building with AI:} In my PolicyPilot project, I built a Gemini based RAG pipeline with citation backed answers to limit hallucination. In my research, my LLM pipeline achieved 25.7x lower latency than the ViBR baseline.
\end{itemize}

I am available for full time roles after my graduation in December 2026. As a Master's student at the University of Cincinnati and someone who values professionalism and growth, I am enthusiastic about contributing to \textbf{Abnormal AI} and helping build secure, cost efficient infrastructure for the \textbf{Platform and Infrastructure team}.

Please find my CV attached below. I look forward to hearing from you. \textbf{I would be excited to help build the platform behind Abnormal AI's security products}.

{\setlength{\parskip}{0pt}
Regards,\\
Tanmay Bhuskute,\\
tanmay.v.bhuskute@gmail.com $\mid$ +1 (732) 858-3027,\\
Cincinnati, OH.
}

\end{letter}

\end{document}
